% !TEX TS-program = xelatex
% !TEX encoding = UTF-8 Unicode
% !Mode:: "TeX:UTF-8"
% Tailored for: Biological / Environmental Sciences Consultant roles

\documentclass{my-resume}
\usepackage{linespacing_fix}
\usepackage{cite}

\begin{document}
\pagenumbering{gobble}

\headerNoPhoto{%
  \name{Xiaoqi Wu}
  \contactInfo{+86 13776020628}{yukiwxq@gmail.com}{Currently based in: London / Suzhou}{Biological \& Environmental Sciences Consulting \textbar\ Shanghai / Suzhou / Hangzhou}
}

\section{Education}
\entry{Imperial College London}{2025.09 - 2026.09}{MSc Computational Methods in Ecology and Evolution}{London}
\entry{Imperial College London}{2022.09 - 2025.07}{BSc Biological Sciences}{London}
\begin{itemize}[parsep=0.2ex]
  \item \textbf{Relevant Coursework}: Biodiversity and Conservation Biology, Global Change Biology, Ecology and Evolution, Behavioural Ecology, Field Skills in Ecology, Fundamentals of Ecology, Field Course in African Biology, Evolution and Diversity
\end{itemize}

\section{Internship \& Engagement}
\entry{Shanghai Sunherb Pharmaceutical Technology Co., Ltd.}{2024.07 - 2024.09}{Analytical Engineer \textperiodcentered\ Analytical Department}{Shanghai}
\begin{itemize}[parsep=0.2ex]
  \item \textbf{Decision-Ready Analysis}: Conducted stability testing of client-supplied drug samples using HPLC/UV, systematically identifying impurity risk through rigorous quality-control procedures to provide clients with analytical data usable directly for stability-assessment decisions
  \item \textbf{Process \& Traceability}: Independently ran the full sample pre-treatment, instrument operation, and data-logging workflow, ensuring accuracy and traceability of results as a reliable basis for cross-departmental quality reporting
\end{itemize}
\entry{Shanghai Jianteng Education Technology Co., Ltd.}{2023.07 - 2023.08}{Teaching Assistant \textperiodcentered\ Biochemistry Department}{Shanghai}
\begin{itemize}[parsep=0.2ex]
  \item \textbf{Teaching Delivery}: Provided extracurricular tutoring for students preparing for Oxford, Cambridge, and other top universities, independently delivering part of the course in English, mock interviews, and class discussions
  \item \textbf{Needs Analysis \& Outcome}: Compiled and analyzed student needs documentation to give teachers differentiated instruction recommendations, helping multiple students secure Cambridge undergraduate offers -- demonstrating structured analysis and client-facing communication
\end{itemize}

\section{Research \& Project Experience}

\entry{Review: Drivers, Impacts, and Conservation Strategies for Insect Decline}{2024.03 - 2025.06}{}{London}
\role{Literature review, critical analysis}{}
\begin{itemize}[parsep=0.2ex]
  \item \textbf{Evidence Synthesis}: Systematically reviewed 70+ papers analyzing drivers (habitat loss, pesticide use, climate change), including a global meta-analysis reporting ~2.5\%/year insect biomass decline
  \item \textbf{Impact Assessment}: Assessed the impact of insect decline on ecosystem services (pollination, nutrient cycling), agriculture, and human health, producing a policy-relevant review
  \item \textbf{Critical Analysis}: Critically evaluated conflicting evidence, a "winners and losers" counter-narrative, and geographic\slash taxonomic reporting bias, proposing future research and policy directions
\end{itemize}

\entry{Drivers of Change in North Atlantic Fisheries Data and Food Web Construction}{2025.01 - 2025.02}{}{London}
\role{Excel, literature review, food web analysis}{}
\begin{itemize}[parsep=0.2ex]
  \item \textbf{Data Analysis}: Quantified a 90-year time series across 8 fish species, identifying species vulnerability such as sole's ~40\% decline (1990-2000) from bycatch and slow maturity
  \item \textbf{Recommendation}: Synthesized policy/climate/fishing-activity literature and proposed the common guillemot as an indicator species to inform fisheries resource management -- report received high marks and distinction
\end{itemize}

\entry{Effect of Environmental Flower Colour Ratios on Insect Colour Preference}{2024.09 - 2024.10}{}{Cape Town}
\role{Field experiment, R (t-test, Pearson correlation)}{}
\begin{itemize}[parsep=0.2ex]
  \item \textbf{Experimental Design}: Designed a 5-colour pan-trap field experiment across 6 sites during a South Africa field course to study insect flower-visiting colour preferences relative to environmental flower colour ratios
  \item \textbf{Results}: Yellow, blue, and white traps all significantly attracted insects over baseline (p$<$0.001); yellow capture rate was positively correlated with the proportion of yellow flowers in the environment (r=0.369, p=0.027), supporting the hypothesis that insects use common environmental flower colours as foraging cues
\end{itemize}

\section{Skills \& Other}
\entrySkillsSep
\begin{itemize}[parsep=0.2ex]
  \item \textbf{Ecology \& Fieldwork}: field experimental design (e.g. multi-site pan-trap colour-preference experiments), behavioural assays with inter-observer reliability testing, biodiversity/environmental surveying, stakeholder engagement and reporting, macroecology and phylogenetic comparative analysis (AVONET-scale trait data, PGLS)
  \item \textbf{Scientific Writing \& Analysis}: literature synthesis and critical analysis, policy-relevant recommendation writing, R/Python statistical support (GLM, PGLS, Wilcoxon, t-tests)
  \item \textbf{Languages}: English (IELTS 7.5), Mandarin (native), French (A2) \textbar\ \textbf{Interests}: piano, photography, skiing
\end{itemize}

\end{document}
